\section{Introduction}
\label{sec:introduction}

Surveillance is a powerful tool, used in various public settings to ensure
safety of the public on a daily basis. CCTV, closed circuit television, is
employed widely with operators having to comb through hours of footage everyday
to uphold security. This is a tedious task, and human error is inevitable.
Posing the question, can we automate this process?

One of the many threats to public safety is unattended luggage, this is a
common occurence which, at best, causes inconvenience, and at worst, could be
placed for a more sinister purpose. Numerous incidents at transportation hubs
and public spaces have demonstrated the potential dangers of unattended luggage
and bags, with counter terrorism authorities advising the public and security
personnel to remain vigilant and report any suspicious
items~\citep{protectuk2022unattended}. Historical acts of terrorism, such as
the Boston Marathon bombing, have emphasised the importance of detecting
suspicious items in crowded spaces, highlighting the need for reliable
detection mechanisms in these scenarios.

Compared to other threats, unattended luggage is distinct in its ability to be
discrete. To the human eye, with low resolution CCTV footage, a security threat
may be indistinguishable from legitimate items, presenting a unique challenge
compared to more obvious security threats such as disorderly conduct. False
positives can increase the burden on teams responsible for monitoring, being
too relaxed could have dire consequences. Crucially, the definition of
'suspicious' is inherently temporal. Unlike detecting a weapon, where the
object itself is the threat, a piece of luggage is neutral until it loses its
association with an owner for a specific duration. This introduces a complex
requirement for a system: it must not only detect objects but track ownership
and measure time, maintaining state across multiple frames in a dynamic
environment.

Systems recording footage at all hours of the day generate an overwhelming
amount of data, the sheer volume posing a significant challenge for operators
to review on a regular basis. This implies a costly requirement for large teams
of personnel manually monitoring the footage, diverting resources from other
facets of security, increasing operational expenses, and potentially leading to
fatigue-induced oversights.

Automation in security is already a widely incorporated concept accross the
public sector, with systems such as facial recognition and intrusion detection
providing results that manual monitoring alone cannot achieve; at least without
significant human-resource investment. An automated approach to detecting
suspicious luggage could alleviate the burden on operators, allowing them to
dedicate more attention to pressing matters; require smaller teams to monitor
larger areas, saving resources; and reduce human error, bolstering overall
security.

Edge processing using embedded systems has become an increasingly viable, cost
effective, and practical solution for surveillance applications. Versus a cloud
solution, edge processing offers lower latency, by way of being situated
physically closer to the data source; improved privacy, sensitive data remains
on site and the onus is on the operator to secure it --- not a cloud provider;
and network independence, there is no reliance on a stable internet connection
just the local network~\citep{pamadi2025edgevscloud}. The proliferation of
affordable, high-performance embedded hardware, with enhanced software support,
has made a strong case for edge processing in surveillance; the benefits of
cloud computing's scalability and power are minimised by new embedded
solutions.

Advances in computer vision technology and deep learning have made the issue of
embedded suspicious luggage detection more tractable. Similar automation tasks,
are being successfully implemented in real-world applications.

Therefore, this project aims to design and implement an embedded system for
suspicious luggage detection. The system will bridge the gap between academic
research and practical deployment, focusing on real-time performance on edge
devices. The target for the project is to create a deployable solution that can
integrate with surveillance infrastructure, utilising an NVIDIA Jetson Orin
Nano as a cost-effective piece of hardware suitable for the task.

